\section{Introdução}

\begin{frame}{Relevância Industrial}
  \begin{columns}
    \column{0.5\textwidth}
    \begin{itemize}
      \item Sistemas reator-separador com reciclo são amplamente estudados, com
        aplicações que vão da produção de polietileno \cite{hafele_2005} a reações
        exotérmicas e endotérmicas acopladas \cite{altimari_2008}.
      \item O reciclo do reagente não convertido aumenta a conversão global,
        melhora a seletividade e reduz etapas de purificação \cite{sainio_2020}.
    \end{itemize}
    \column{0.5\textwidth}
    \begin{figure}
      \centering
      \includegraphics[width=\textwidth]{figuras/indorama.jpg}
      \caption{Unidade de Etoxilação. Fonte: Autores, 2026.}
    \end{figure}
  \end{columns}
\end{frame}

\begin{frame}{Motivação e Objetivo}
  \small

  \begin{block}{Motivação}
    \begin{itemize}
      \item O reciclo melhora o desempenho de conversão, mas também torna a análise dinâmica mais desafiadora.

      \item O retorno de material aumenta o acoplamento entre variáveis e intensifica o comportamento não linear do processo \cite{zeyer_2004,sainio_2020,bildea_2000}.

      \item A resposta do sistema pode ser sensível à estratégia operacional adotada \cite{pushpavanam_2001}.
    \end{itemize}
  \end{block}

  \vspace{0.25cm}

  \begin{block}{Objetivo do trabalho}
    Caracterizar o comportamento dinâmico de um sistema reator-separador acoplado por reciclo por meio de ferramentas clássicas de análise de sistemas lineares nos domínios do tempo contínuo, da frequência e do tempo discreto.
  \end{block}

\end{frame}
